% ACL conference style.
% Template: https://www.overleaf.com/latex/templates/
%           association-for-computational-linguistics-acl-conference/jvxskxpnznfj
% acl.sty and acl_natbib.bst are ACL-distributed; see README.md.
%
% Every number in this paper is produced by scripts/build_paper_tables.py
% from the per-run records. Nothing here is typed by hand.
\documentclass[11pt]{article}

\usepackage[review]{acl}

\usepackage{times}
\usepackage{latexsym}
\usepackage[T1]{fontenc}
\usepackage[utf8]{inputenc}
\usepackage{microtype}
\usepackage{graphicx}
\usepackage{booktabs}
\usepackage{multirow}
\usepackage{amsmath}
\usepackage{amssymb}
\usepackage{url}

\title{Can Five Seeds Tell Embedding Initializers Apart?\\
A Controlled Study of Vocabulary Extension for Amharic}

\author{Hailay Kidu Teklehaymanot \\
  Mekelle University, Ethiopia \\
  \texttt{hailay.kidu@mu.edu.et} \\\And
  Dren Fazlija \\
  L3S Research Center \\
  Hannover, Germany \\
  \texttt{dren.fazlija@L3S.de} \\\And
  Wolfgang Nejdl \\
  L3S Research Center \\
  Hannover, Germany \\
  \texttt{nejdl@L3S.de} \\}

\begin{document}
\maketitle

\begin{abstract}
Extending a pretrained encoder's vocabulary to a new language requires an
embedding for every added token, and several initialization strategies have
been proposed to supply one. We ask whether the benchmarks commonly used for
Amharic can distinguish between them. Under a single fixed protocol---one
configuration, one set of rebuilt data splits, one untuned fine-tuning
recipe, shared seeds---we compare morpheme-grounded initialization,
overlap-based initialization, and no vocabulary extension at all, across
text classification, named entity recognition and extractive question
answering. Across six paired comparisons, every 95\% confidence interval
contains zero. On all three tasks the seed-to-seed variation exceeds the
entire spread between the three systems, and on text classification the
smallest difference five seeds could detect at conventional power is 19.8
macro-F1 points against an observed between-system spread of 9.0. A further property of the setup
bounds what any such comparison can show: 6 of 60 runs never fit their
training data and settle into degenerate predictors, distributed across all
three systems and identifiable from training loss alone. We also record that
a favourable comparison reported under a superseded protocol reverses under
this one, and that the two protocols differ in test set, training set and
class prior together, so the shift cannot be attributed to any single cause.
We report means and intervals and claim no ordering. This paper replaces a
retracted predecessor \citep{teklehaymanot-etal-2026-lgse}; all code, data
manifests, per-run records and the script that generates every number are
released.
\end{abstract}

\section{Introduction}

A multilingual encoder adapted to a new language by vocabulary extension
needs a vector for each token it has never seen. Random initialization
discards what is known about the token's form. FOCUS
\citep{dobler-de-melo-2023-focus} builds each new embedding from tokens it
overlaps with in an auxiliary space. A morpheme-grounded alternative
decomposes the token and averages over its morphemes, which is appealing for
Amharic, where affixation is heavy and a frequency-driven tokenizer
fragments surface forms arbitrarily.

Papers in this area report differences of one to several points between such
strategies, with standard deviations below one point, and conclude in favour
of one of them. This paper asks a prior question: at the sample sizes
conventionally used, can these benchmarks support such a conclusion at all?

We answer it for Amharic by running three systems under one protocol and
examining the variance structure rather than only the means. The answer is
that they cannot. We also measure a fine-tuning instability that costs 10\%
of runs, and we document a favourable comparison reported under a superseded
protocol that reverses under this one.

\paragraph{Relation to a retracted paper.} An earlier paper by these authors
described the morpheme-grounded method and reported a results table for which
no archived experimental record exists, a projection matrix that was never
specified, and a dataset that was not constructed. That paper
\citep{teklehaymanot-etal-2026-lgse} is retracted. The present paper replaces
it. It shares no text and no results with it. All data splits were rebuilt
from their original sources and are documented in \S\ref{sec:setup} and
Appendix~\ref{app:prov}, which maps every number reported here to the job
identifier, seed and result file it came from.

\paragraph{Contributions.}
\begin{enumerate}
\itemsep0em
\item A controlled three-way comparison on Amharic---60 runs, one
  configuration, shared seeds---in which no paired confidence interval
  excludes zero (\S\ref{sec:results}).
\item A measurement of why: seed variation exceeds between-system spread on
  every task, and the minimum detectable difference at five seeds is an
  order of magnitude larger than the effects under discussion
  (\S\ref{sec:resolution}).
\item A characterization of a fine-tuning failure mode affecting 6 of 60
  runs, separable from the method under test using training loss alone
  (\S\ref{sec:failures}).
\end{enumerate}

\section{Systems}

\paragraph{XLM-R.} Off-the-shelf XLM-R base
\citep{conneau-etal-2020-unsupervised}, fine-tuned directly on each task. No
vocabulary extension, no adaptive pretraining, no auxiliary embedding space.
This is the far control: it bounds what the rest of the machinery is worth.

\paragraph{FOCUS.} Vocabulary extension with FOCUS initialization
\citep{dobler-de-melo-2023-focus}, followed by language-adaptive pretraining.
Each new embedding is a weighted combination of overlapping tokens in a
FastText space, projected into the model's space. This is the near control:
it consumes the same auxiliary space and the same projection as the
morpheme-grounded system, and differs only in not using morpheme
segmentation.

\paragraph{LGSE.} The morpheme-grounded initializer. A new token $t$ with
decomposition $m_1 \dots m_k$ receives
\begin{equation}
e_{\text{new}} = \frac{1}{k}\sum_{j=1}^{k} W e_{m_j},
\label{eq:lgse}
\end{equation}
where $e_{m_j}$ is a FastText vector and $W \in \mathbb{R}^{d \times d}$
aligns the two spaces. Tokens without a decomposition fall back to a
whole-token vector and then to character $n$-grams. During adaptive
pretraining a term $\lambda \lVert e_{\text{new}} - \mu \rVert^2$ penalizes
drift from the initialization.

\paragraph{Two quantities the method does not fix.} Equation~\ref{eq:lgse}
requires $W$, and the regularizer requires $\lambda$. Neither is determined
by the method as described, and the retracted paper specified neither. We
therefore chose both, documented the choice, and held it fixed for every
run: $W$ is obtained by orthogonal Procrustes,
$W^\star = \arg\min_{W^\top W = I}\lVert WX - Y\rVert_F$, on 5{,}910 surface
anchors shared between the Amharic FastText model and XLM-R (residual
Frobenius norm 423.0 against 633.0 unaligned), and $\lambda = 1.0$. These are
our choices. A reader should treat the LGSE column as one faithful
instantiation rather than the only one.

We note one structural consequence. The initializer writes $We_{m_j}$ into
the embedding matrix and the regularizer anchors to a constant, so nothing
downstream is a function of $W$ once initialization is complete. $W$ is a
fixed preprocessing transform here, not a learned parameter.

\section{Setup}
\label{sec:setup}

\paragraph{One protocol.} All 60 runs share a single configuration file,
verified by checksum: every run records the same \texttt{config\_sha256}.
Fine-tuning uses a constant learning rate of $5\times10^{-5}$, no warmup,
batch size 32, 10 epochs, and checkpoint selection by evaluation loss. This
recipe is identical for all three systems and was not tuned for any of them;
the selection metric is task-agnostic by construction. Appendix
\ref{app:prov} records the checksum and the full recipe.

\paragraph{Data.} For NER we take the MasakhaNER splits
\citep{adelani-etal-2021-masakhaner} unchanged, and for question answering
the AmQA splits \citep{taffa-etal-2024-amqa} unchanged. Text classification
uses a
2{,}499-row annotated educational-quality corpus, reduced to 2{,}340 after
removing 2 spreadsheet-error rows and 157 exact duplicates, split
1872/233/235 by a stratified procedure at seed 42 and verified
leakage-free. The first author led that annotation. Each split carries a
manifest with source, raw checksum and split seed.

\paragraph{Metrics.} F1 for NER and QA. For text classification we report
macro-F1 over the five classes present in the test split as the primary
metric, with accuracy secondary. The majority class holds 66.4\% of the test
set, so accuracy is dominated by it: the majority-class baseline is 66.38
accuracy against 15.96 macro-F1. We also report the macro-F1 value recorded
in the run files, which averages over six classes including one with no test
instances and is therefore capped at five sixths of the attainable maximum.

\paragraph{Seeds.} 42--46 for every cell, extended to 42--51 on NER for all
three systems after the failure mode of \S\ref{sec:failures} appeared there.

\paragraph{Scope.} Amharic only, three systems. Two further systems from the
retracted design---adaptive pretraining alone, and extension with random
initialization---were not run under this protocol, and
\S\ref{sec:limitations} explains why results from an earlier batch cannot
substitute for them.

\section{Results}
\label{sec:results}

\begin{table}[t]
\centering
\small
\begin{tabular}{llccc}
\toprule
Task & & LGSE & FOCUS & XLM-R \\
\midrule
\multirow{4}{*}{TC}
 & macro-F1 & 45.61 & 54.56 & 45.59 \\
 & sd       & 12.80 & 2.59  & 15.04 \\
 & accuracy & 78.30 & 81.62 & 78.98 \\
 & $n$ / fail & 5 / 2 & 5 / 0 & 5 / 2 \\
\midrule
\multirow{3}{*}{NER}
 & F1  & 61.44 & 60.10 & 67.05 \\
 & sd  & 21.64 & 21.30 & 2.55 \\
 & $n$ / fail & 10 / 1 & 10 / 1 & 10 / 0 \\
\midrule
\multirow{3}{*}{QA}
 & F1  & 58.23 & 59.41 & 59.54 \\
 & sd  & 1.40  & 1.51  & 2.14 \\
 & $n$ / fail & 5 / 0 & 5 / 0 & 5 / 0 \\
\bottomrule
\end{tabular}
\caption{All 60 runs, one protocol. TC macro-F1 is over the five classes
present in test; the majority-class baseline is 15.96 macro-F1 and 66.38
accuracy. \emph{fail} counts runs that never fit their training data
(\S\ref{sec:failures}); they are retained here. No ordering is claimed:
see Table~\ref{tab:paired}.}
\label{tab:main}
\end{table}

\begin{table}[t]
\centering
\small
\begin{tabular}{llcrc}
\toprule
Task & vs & $n$ & $\Delta$ & 95\% CI \\
\midrule
\multirow{2}{*}{TC}  & FOCUS & 5  & $-8.95$ & $[-23.74, +5.84]$ \\
                     & XLM-R & 5  & $+0.03$ & $[-20.53, +20.58]$ \\
\midrule
\multirow{2}{*}{NER} & FOCUS & 10 & $+1.34$ & $[-21.61, +24.29]$ \\
                     & XLM-R & 10 & $-5.61$ & $[-22.03, +10.80]$ \\
\midrule
\multirow{2}{*}{QA}  & FOCUS & 5  & $-1.17$ & $[-2.91, +0.56]$ \\
                     & XLM-R & 5  & $-1.30$ & $[-3.12, +0.51]$ \\
\bottomrule
\end{tabular}
\caption{Paired differences, LGSE minus comparator, on matched seeds.
Positive favours LGSE. \textbf{Every interval contains zero.} Exact
Wilcoxon agrees (0.188--0.812); at these sample sizes the smallest
attainable two-sided exact $p$ is 0.062 at $n=5$ and 0.125 at $n=4$, so four
of these six rows could not have reached $p<0.05$ at any effect size, and we
do not report $p$-values as evidence.}
\label{tab:paired}
\end{table}

Table~\ref{tab:main} gives the three systems under the common protocol and
Table~\ref{tab:paired} the paired differences.

\paragraph{No interval excludes zero.} The largest observed difference is
$-8.95$ macro-F1 on text classification against FOCUS, with an interval
spanning $-23.74$ to $+5.84$. The smallest is $+0.03$ against XLM-R, an
interval 41 points wide. Nothing in Table~\ref{tab:paired} supports an
ordering of the three systems.

\paragraph{Win counts are more informative than the means, and still
insufficient.} LGSE loses 4 of 5 QA seeds to each comparator and wins 0 of 5
text-classification seeds against FOCUS, which is a consistent direction.
But five paired observations cannot certify a direction at conventional
thresholds, and \S\ref{sec:resolution} shows the margin by which they fall
short.

\paragraph{Off-the-shelf XLM-R is not separated from either extended
system.} It reaches the highest mean on NER and QA while performing no
vocabulary extension and no adaptive pretraining, at a fraction of the
compute: its runs took 6--24 minutes against roughly 5 hours for the
extended systems. We state this as a non-separation rather than a result in
its favour; its text-classification standard deviation is the largest in the
table, at 15.04.

\section{Why the Comparison Does Not Resolve}
\label{sec:resolution}

\begin{table}[t]
\centering
\small
\begin{tabular}{lrrr}
\toprule
Task & Spread & Mean sd & Ratio \\
\midrule
TC  & 8.98 & 10.14 & 0.89 \\
NER & 6.95 & 15.16 & 0.46 \\
QA  & 1.30 & 1.68  & 0.77 \\
\bottomrule
\end{tabular}
\caption{Between-system spread against within-system seed variation, all
runs. A ratio below one means changing the seed moves a system more than
changing the system does. All three are below one.}
\label{tab:snr}
\end{table}

Table~\ref{tab:snr} compares the spread between the three systems against
the spread between seeds of one system. On every task the second exceeds the
first. Excluding the failed runs of \S\ref{sec:failures} does not change the
conclusion: the ratios become 0.56, 0.64 and 0.77.

The same point in units of detectable effect: using the observed
paired-difference standard deviations, the minimum difference detectable at
$\alpha = 0.05$ and 80\% power is 19.80 macro-F1 on text classification, 3.10
F1 on NER and 2.32 F1 on QA. The text-classification figure is more than
twice the entire observed spread between systems. Reaching a 2-point
detectable difference there would require approximately 281 seeds per
system. We report this as a property of the benchmark and the sample size,
not of any system.

\section{A Comparison That Did Not Survive Re-measurement}
\label{sec:split}

An earlier batch of runs, under a superseded protocol, reported the
morpheme-grounded system ahead of FOCUS on text classification by 1.97
accuracy points. Under the protocol of \S\ref{sec:setup} it trails by 3.32
and wins no seed. We record this because the direction of a published
comparison reversed, and because the reason is instructive about comparing
across protocol revisions.

\begin{table}[t]
\centering
\small
\begin{tabular}{lcc}
\toprule
 & Superseded & Current \\
\midrule
Test items            & 71     & 235 \\
Majority-class share  & 35.21\% & 66.38\% \\
Documents in corpus   & 700    & 2{,}340 \\
\midrule
LGSE (accuracy)       & 75.77  & 78.30 \\
FOCUS (accuracy)      & 73.80  & 81.62 \\
LGSE $-$ FOCUS        & $+1.97$ & $-3.32$ \\
\bottomrule
\end{tabular}
\caption{The two text-classification protocols. They differ in test set,
training set and class prior together, so the accuracy columns are not two
measurements of one quantity and the per-system shifts are not interpretable
as effects. What the table shows is that the comparison between the two
systems reverses. Per-class supports are in Appendix~\ref{app:tcsplit}.}
\label{tab:split}
\end{table}

\paragraph{The two protocols are not two measurements of the same task.}
It would be convenient to read FOCUS's movement from 73.80 to 81.62 as the
effect of a larger test set on a system that performs no morpheme
segmentation. That reading is not available. Rebuilding the corpus changed
the training split and the label distribution as well as the test split: the
majority class holds 35.21\% of the superseded test set and 66.38\% of the
current one, so a classifier that predicts only the majority label gains
31.17 accuracy points from the change alone, with no change in capability.
Accuracy on the two splits is scaled differently, and no per-system
difference between the columns can be attributed to a single cause.

\paragraph{What the reversal does establish.} Within the current protocol,
measured on one split with one configuration and shared seeds, the
morpheme-grounded system does not lead FOCUS on text classification; it
trails by 3.32 points and loses every seed. That comparison is internally
valid and is the one reported in Table~\ref{tab:main}. The earlier
favourable comparison is superseded, not contradicted by a controlled
intervention, and we do not claim to have isolated why it differed.

\paragraph{Consequence for the missing systems.} Two systems from the
original design were not run under the current protocol
(\S\ref{sec:limitations}). Their numbers exist only in the superseded batch,
on the 71-item split with a 35.21\% majority class. Placing those values in
Table~\ref{tab:main} beside systems measured against a 66.38\% majority
class would produce an ordering driven by the class prior rather than by the
systems. We therefore leave those rows empty rather than fill them across
protocols.

\section{A Fine-Tuning Failure Mode}
\label{sec:failures}

Six of 60 runs never fit their training data. Their relative
training-loss reduction from the first to the last logged epoch falls between
21.7\% and 28.9\%, against 69.5\% to 97.9\% for every other run---a gap with
no runs inside it. Any threshold between 28.9\% and 69.5\% selects the same
six, so the classification does not depend on where in the gap it is placed.
The rule uses training loss only and never a test score.

\begin{table}[t]
\centering
\small
\begin{tabular}{llcc}
\toprule
Task & System & Seed & Loss reduction \\
\midrule
NER & FOCUS & 44 & 24.1\% \\
NER & LGSE  & 45 & 21.7\% \\
TC  & LGSE  & 42 & 28.9\% \\
TC  & LGSE  & 43 & 24.8\% \\
TC  & XLM-R & 42 & 21.9\% \\
TC  & XLM-R & 45 & 26.8\% \\
\bottomrule
\end{tabular}
\caption{The six runs selected by the training-loss rule. They occur under
all three systems, so the instability belongs to the fine-tuning setup
rather than to any initializer.}
\label{tab:fail}
\end{table}

The resulting predictors are degenerate in a task-specific way. On NER the
tagger proposes no entity at all, giving precision, recall and F1 of exactly
zero. On text classification the classifier never predicts class 3, which
holds 19\% of the test set, and never predicts the two rarest classes;
per-class F1 shows this as a partial label collapse rather than a total one.
Dev scores track test scores in each case, so the failure is visible without
touching the test set.

These runs are not checkpoint mis-selection: the evaluation-loss minimum
across all ten saved checkpoints is 0.77--0.80 in the failed runs against
0.42--0.55 in the successful ones, so no available checkpoint was good. Nor
is it overfitting: training loss stays high.

We retain all six runs in Table~\ref{tab:main} and report a secondary
analysis excluding them. Excluding them moves the text-classification means
to 54.67, 54.56 and 56.45 for LGSE, FOCUS and XLM-R, on three, five and
three runs, and the NER means to 68.27, 66.78 and 67.05 on nine, nine and
ten runs. Neither policy separates the systems: the intervals remain wide
and overlapping under both. We report both because the choice moves the
text-classification means by nine points and the NER means by seven, and we
see no principled basis for preferring one.

\section{Recommendations}

Each observation above has a cheap remedy that we would have benefited from.

\paragraph{Check resolution before comparing.} Whether the between-system
spread exceeds the between-seed spread is computable from a pilot, before any
method is proposed. Where it does not, as in three of three tasks here, a
five-seed comparison cannot support a ranking regardless of which system
wins.

\paragraph{Report the minimum detectable difference.} One line per
comparison. It distinguishes ``no difference was found'' from ``no difference
of this size could have been found''.

\paragraph{Run a system that shares everything but the contribution.} FOCUS
differs from LGSE only in morpheme segmentation, so the pair bounds what the
segmentation itself can be worth. A system compared only against an
unmodified baseline cannot separate the contribution from everything else
that changed alongside it.

\paragraph{Do not compare across protocol revisions.} When a corpus is
rebuilt, the test set, the training set and the class prior can all move
together, as they did here (\S\ref{sec:split}): the majority-class baseline
went from 35.21\% to 66.38\%, which alone rescales accuracy by 31 points.
Numbers from before and after such a change are not two measurements of one
quantity, and a table that mixes them reports the revision rather than the
systems. Re-run the baselines or leave the cells empty.

\paragraph{Release the operator, not only the equation.} $W$ in
Equation~\ref{eq:lgse} is 2.3\,MB. Releasing it, or the code that builds it,
removes the largest obstacle to reproducing a method of this kind. We release
ours with a sidecar recording its language, anchor count, objective and
residual.

\paragraph{Disclose degenerate runs.} A single all-\texttt{O} seed can move a
mean by 13 points and inflate a standard deviation past 30. Neither silent
exclusion nor silent averaging is adequate; report both and state the rule.

\section{Related Work}

\paragraph{Initialization for vocabulary extension.} WECHSEL
\citep{minixhofer-etal-2022-wechsel} transfers embeddings through bilingual
dictionaries and static vectors. FOCUS \citep{dobler-de-melo-2023-focus},
our near control, builds each new embedding from overlapping tokens in an
auxiliary FastText space. Both FOCUS and the morpheme-grounded method reduce
to a whole-token vector when their respective condition---overlap, or a
decomposition---is unavailable, which is why they are the right pair to
compare and why their difference is bounded by how often that condition
fires. \citet{csaki-etal-2023-efficiently} survey the adaptation pipeline.
Our results do not argue against this line of work; they indicate that these
benchmarks at these sample sizes cannot adjudicate between its members.

\paragraph{Morphological resources for Ethio-Semitic languages.} Amharic is
templatic and heavily affixing, which motivates morpheme-aware vocabulary
construction. HornMorpho \citep{gasser-2011-hornmorpho} provides finite-state
analysis for Amharic and Tigrinya and is enabled in the configuration used
here; with it, 99 of the 198 added Amharic tokens (50.0\%) receive a
decomposition into more than one morpheme, against 33 (16.7\%) from the
static lexicon alone. We report these as properties of the pipeline and do
not use them to explain the comparison outcomes.

\paragraph{Reproduction and statistical practice.}
\citet{belz-etal-2021-systematic} survey reproducibility in NLP. Our
methodology is control-based: rather than reporting that numbers differ, we
hold systems that should be unaffected by a change fixed and attribute the
difference from their movement.

\section{Conclusion}

Under one protocol, on three Amharic tasks, no paired comparison between
morpheme-grounded initialization, overlap-based initialization and no
extension at all excludes zero. The reason is measurable and specific:
changing the seed moves a system further than changing the system does, the
smallest detectable difference on text classification is twice the observed
spread, and a tenth of runs fail to fit their training data. A comparison
reported under a superseded protocol reverses under this one.

We do not conclude that these initializers are equivalent; our tests are too
small to support that, and we say so. We conclude that this evaluation suite,
at the sample sizes the field conventionally uses, does not answer the
question it is being asked, and that the four properties above are worth
measuring before the next comparison of this kind is published.

\section*{Limitations}
\label{sec:limitations}

Our $W$ and $\lambda$ are documented substitutions, not values recovered
from the original method, so the LGSE column is one faithful instantiation
and the absolute numbers would differ under another.

Two systems from the original design---adaptive pretraining alone, and
extension with random initialization---were not run under this protocol. An
earlier batch contains them, but evaluated text classification on the
superseded 71-item test set, carries three different configuration
checksums, and mixes two alignment matrices; \S\ref{sec:split} shows that
placing those numbers beside these would create an ordering artifact of
about 8 points. We therefore leave the rows empty rather than fill them from
a different protocol. Completing them under this protocol would add two rows
at five seeds each, which Table~\ref{tab:snr} and the detectable-difference
figures indicate would not change the conclusion.

The paired tests have five to ten observations. They establish that these
data do not separate the systems, not that the systems are equivalent;
distinguishing those readings would need many more seeds or benchmarks with
a larger between-system spread.

This study covers Amharic. No cell for a second language exists under this
protocol, so we make no cross-linguistic claim.

The text-classification corpus is small, its label distribution is heavily
skewed, and one of its six labels has no test instances. Our primary metric
excludes that label; the recorded six-class value is reported alongside for
traceability.

\section*{Ethics Statement}

This paper replaces a retracted paper by the same authors. We report the
outcome as measured, including where it contradicts our own earlier claims,
and we document errors made during this work alongside those found in the
original: a configuration error that applied the Tigrinya alignment matrix
to Amharic runs, a morpheme-coverage figure initially over-reported before
re-measurement, and three log-parsing defects in our own analysis code, each
of which produced a wrong failure classification and is now pinned by a
regression test. We found no evidence of fabrication in the original work;
every discrepancy we traced is consistent with ordinary pipeline defects of
the kind we also made here.

Self-examination carries a conflict of interest. We mitigate it by
pre-registering the one directional control before it ran
(\S\ref{sec:split}), by releasing every per-run record so that each number
can be recomputed independently, and by reporting the outcome least
favourable to our own prior claim. Readers should weigh our conclusions
accordingly.

All datasets are publicly available or released by us. The Amharic
educational-quality annotation was led by the first author. No human
subjects were involved.

\section*{Acknowledgments}

We thank the reviewers of the earlier submission, whose questions prompted
this study, and the L3S compute support team.

\bibliography{custom}

\appendix

\section{Provenance}
\label{app:prov}

\paragraph{One script.} Every number in this paper is produced by
\texttt{scripts/build\_paper\_tables.py}, which reads the per-run
\texttt{experiment.json} records and emits Tables \ref{tab:main},
\ref{tab:paired}, \ref{tab:snr}, \ref{tab:fail}, the per-class
text-classification breakdown, the coverage counts, and a per-seed CSV.
Running it reproduces every table; nothing in this paper was typed by hand.

\paragraph{Runs.} 60 runs, all from
\texttt{results\_hornmorpho\_wam/}: three systems $\times$ three tasks
$\times$ five seeds, with NER extended to ten. Each record carries its
commit, configuration checksum, alignment-matrix path, corpus checksum and
per-seed dev and test metrics. All 60 share
\texttt{config\_sha256} \texttt{3e51cdab\ldots}.

\paragraph{Excluded directories.} Four other result directories exist and are
excluded for protocol differences, never for their outcomes:
\texttt{results\_gen2\_snapshot\_20260924} and \texttt{results} (150 runs
each; three configuration checksums, two alignment matrices, 71-item
text-classification test set), \texttt{results\_hornmorpho} (54 runs; mixes
the two alignment matrices), and \texttt{results\_tigrinya\_gen2} (45 runs;
Tigrinya only, no three-system comparison).

\paragraph{Recipe.} Adaptive pretraining: frozen backbone, maximum length
256, MLM probability 0.15, learning rate $5\times10^{-5}$ constant, batch
size 32, 10 epochs, weight decay 0.01, AdamW, fp16. Fine-tuning: the same
learning rate, schedule, batch size and epoch count, with checkpoint
selection by evaluation loss. Identical for all three systems, untuned.
FastText: skipgram, dimension 768, 5 epochs, minimum count 5, vocabulary
165{,}471. Vocabulary extension adds 198 Amharic tokens.

\paragraph{Compute.} Each extended-system run takes approximately 5 hours on
one A100, dominated by adaptive pretraining, which is repeated per seed
because the seed determines the initialization under test. XLM-R runs take
6--24 minutes. The 60 runs total roughly 230 GPU-hours.

\paragraph{Per-seed values.} Released as \texttt{tables/per\_seed.csv}: one
row per (task, system, seed) with the score, the failure flag, the commit and
the result file path.

\section{The Two Text-Classification Protocols}
\label{app:tcsplit}

Per-class test supports, read from the run records of both batches. The
superseded protocol is retained here only to document why its accuracy
figures are not comparable with the current ones (\S\ref{sec:split}); no
number from it enters Table~\ref{tab:main}.

\begin{table}[h]
\centering
\small
\begin{tabular}{lrrrrrrr}
\toprule
Label & 1 & 2 & 3 & 4 & 5 & 6 & total \\
\midrule
Superseded & 25 & 24 & 13 & 7 & 0 & 2 & 71 \\
Current    & 25 & 156 & 45 & 7 & 0 & 2 & 235 \\
\bottomrule
\end{tabular}
\caption{Class 2 grows from 24 to 156 instances, taking the majority share
from 35.21\% to 66.38\%. Label 5 occurs in training but in neither
development nor test split under either protocol, which is why the recorded
six-class macro-F1 is capped at five sixths of its maximum
(\S\ref{sec:setup}).}
\end{table}

\end{document}
